\section{文献综述}
\label{sec:lit-rev}

本节从三个方面综述相关工作：自动化文献综述系统、研究任务的多Agent架构、以及LLM文本生成中的反幻觉技术。

\subsection{自动化文献综述系统}

文献综述过程的自动化日益受到研究关注。早期方法聚焦于基于关键词的检索和文献计量分析 \citep{white-mccain-1998}。近年来，\citet{wang-etal-2020}提出了使用抽象摘要的自动化综述生成框架。\citet{alaofi-etal-2023}探索了使用大语言模型生成文献综述的方法，展示了LLM生成综述的潜力和风险。

\citet{kostousov-2025}在IEEE Access上提出了AI增强的文献综述方法论，结合语义聚类和摘要生成。我们的系统在此方法论基础上引入了多项关键增强：（1）层级RQ驱动的流水线替代扁平主题聚类；（2）GPU加速的嵌入计算；（3）聚焦底层原理的提示词工程，强制要求公式级分析。

\subsection{多Agent架构}

多Agent系统已广泛应用于复杂任务分解。\citet{park-etal-2023}展示了基于LLM的Agent在协作场景中的有效性。\citet{hong-etal-2024}提出了MetaGPT，一个用于多Agent协作的元编程框架。本系统采用的协调者-代理模式受到这些工作启发，但专门针对文献综述领域进行了优化，引入了领域特定的质量门禁和检查点机制。

\subsection{LLM生成中的反幻觉技术}

幻觉——生成事实上不正确或虚构的内容——仍然是LLM应用中的关键挑战 \citep{ji-etal-2023}。在文献综述场景中，幻觉主要表现为虚构引用、错误的方法描述和缺乏依据的论断。\citet{shuster-etal-2021}证明检索增强生成（RAG）可以通过将生成内容锚定在检索到的文档上来减少幻觉。本系统采用多层反幻觉策略：严格的引用格式强制（附正确/错误示例）、禁止空泛表述、强制公式级推理，以及回退报告机制实现优雅降级。
